(1)$\Rightarrow$(2): If $\overline{E\cap X_0}\neq E$, its complement meets $E$ in a nonempty locally closed set disjoint from $X_0$, a contradiction. Conversely, (2) makes $E\cap X_0$ dense in $E$, so it meets every nonempty relatively open subset $U\cap E$.

(1)$\Rightarrow$(3): Restriction is surjective by the subspace topology. If two open sets differ, one of their differences is a nonempty locally closed set, so it meets $X_0$ and their restrictions differ. Conversely, if $U\cap E$ is nonempty and misses $X_0$, then the distinct open sets $U$ and $U\setminus E$ have the same restriction, contradicting (3).

For $X=\operatorname{Spec}A$ and $X_0$ the maximal ideals, the closure of $V(I)\cap X_0$ is $V(\bigcap_{\mathfrak{m}\supseteq I,\ \mathfrak{m}\text{ maximal}}\mathfrak{m})$. It equals $V(I)$ for every $I$ exactly when every quotient has equal Jacobson radical and nilradical. Apply \exref{5.23}.

Suppose $A$ is Jacobson and $\{\mathfrak{p}\}$ is locally closed. There is $f\notin\mathfrak{p}$ with $V(\mathfrak{p})\cap D(f)=\{\mathfrak{p}\}$. If $\mathfrak{p}$ were not maximal, \exref{5.23}(iii) would give a prime strictly containing it and avoiding $f$, a contradiction. Thus the singleton is closed. Conversely, if condition (iii) of \exref{5.23} fails, there are a nonmaximal prime $\mathfrak{p}$ and $f\notin\mathfrak{p}$ lying in every prime strictly containing it. Then $V(\mathfrak{p})\cap D(f)=\{\mathfrak{p}\}$ is locally closed and not closed.
